\begin{frame}{Construire un chemin de localisation}
    \begin{itemize}
        \item Un chemin XPath est une suite d'étapes
        \item Les étapes sont séparées par des "/"
    \end{itemize}
    \textbf{Une étape est constituée de trois composants:}
    \begin{itemize}
        \item \textbf{Axe =} sens de parcours des noeuds (\texttt{parent::}, \texttt{descendant::})
        \item \textbf{Test de noeud =} nom d'un noeud dans le sens de l'axe
        \item \textbf{Prédicat(s) =} précise la caractéristique du noeud (sorte de filtre)
    \end{itemize}
    \texttt{//lg[position()=2]/\textcolor{red}{descendant::}\textcolor{green}{l} \textcolor{blue}{[position()=2]}/text()}
\end{frame}

\begin{frame}{Syntaxe des noeuds}
    \begin{itemize}
        \item \textbf{Élément:} Désigne l'élément.
        \item \textbf{@NomAttribut:} Désigne l'attribut.
        \item \textbf{text():} Désigne le contenu textuel d'un élément sans la balise.
        \item \textbf{comment():} Désigne le contenu d'un commentaire.
    \end{itemize}
\end{frame}

\begin{frame}{Premier composant d'une étape: axe}
        \include{03 - XPath/syntaxe_tableau}    
\end{frame}

\begin{frame}{Deuxième composant: le "test de noeud"}
    \textrightarrow{} Soit le nom du noeud soit le type de noeud
    \begin{itemize}
        \item \textbf{nom du noeud:} Vers le noeud dont le nom est spécifié
        \item \textbf{*:} Vers tous les noeuds
        \item \textbf{node():} Vers tous les types de noeuds (éléments, commentaires, attributs, etc...)
        \item \textbf{text():} Vers les noeuds de type texte
        \item \textbf{comment():} Vers les noeuds de type commentaire
    \end{itemize}
\end{frame}

\begin{frame}{Exemple axe + filtre}
    \begin{itemize}
        \item Je veux récupérer le \texttt{<head>} en chemin absolu:
        \texttt{/\textcolor{red}{child::}\textcolor{blue}{TEI}/\textcolor{red}{child::}\textcolor{blue}{text}/\textcolor{red}{child::}\textcolor{blue}{body}/\textcolor{red}{child::}\textcolor{blue}{head}}
    \end{itemize}
\end{frame}

\begin{frame}{Exemple axe + filtre}
    \begin{figure}
        \centering
        \includegraphics[width=\linewidth]{03 - XPath/arbre_xml_head_path.png}
    \end{figure}
\end{frame}

\begin{frame}{Exemple axe + filtre}
    \begin{itemize}
        \item On peut aussi récupérer le \texttt{<head>} en chemin relatif:\\
        \texttt{//body/child::head}
    \end{itemize}
\end{frame}

\begin{frame}{Exemple axe + filtre}
Récupérer tous les commentaires du document TEI:
    \begin{itemize}
        \item Se fait en une seule étape en sélectionnant tous les descendants du noeud racine qui sont des commentaires:\\
        \texttt{/descendant::comment()}        
    \end{itemize}
\end{frame}

\begin{frame}{Troisième composant: le prédicat}
    \begin{itemize}
        \item On peut filtrer grâce aux prédicats
        \item \texttt{Axe::Noeud[Prédicat1][Prédicat2]...}
    \end{itemize}
    \texttt{//lg[position()=2]/descendant::l \textcolor{blue}{[position()=2]/text()}}
\end{frame}

\begin{frame}{Troisième composant: le prédicat}
    \begin{itemize}
        \item \textbf{attribute:} Affiner la recherche en fonction d'un attribut
        \item \textbf{count():} Compter le nombre de noeuds
        \item \textbf{last():} Sélectionner le dernier noeud d'une liste
        \item \textbf{position():} Affiner la recherche selon la position d'un noeud
    \end{itemize}
\end{frame}

\begin{frame}{Troisième composant: le prédicat}
    \begin{block}{Exemples:}
        \begin{itemize}
            \item \textcolor{blue}{\texttt{//lg[@type]}}
            \begin{itemize}
                \item Les noeuds lg qui ont un attribut type
            \end{itemize}
            \item \textcolor{blue}{\texttt{//lg[@type="quatrain"]}}
            \begin{itemize}
                \item Les noeuds lg qui un attribut type valant quatrain
            \end{itemize}
            \item \textcolor{blue}{\texttt{//lg/descendant::text()[position()=1]}}
            \begin{itemize}
                \item Le premier noeud de type texte descendant d'un \texttt{<lg>}
                \item Peut s'abréger: \textcolor{blue}{ \texttt{//lg/descendant::text()[1]}}
            \end{itemize}
        \end{itemize}
    \end{block}
\end{frame}